% =====================================================================
% Track B (Open-World, pretrained backbones allowed): V-JEPA 2 pipeline.
%
% \input{}-ed from report/main.tex.  Section labels use the \texttt{trackb:*}
% prefix so they do not clash with Track~A \texttt{tracka-phase4:*} labels.
% =====================================================================

% Set Track B accent color (blue stripe + blue headers)
\renewcommand{\trackcolor}{trackblue}
\renewcommand{\trackname}{Track~B}

\section{Track B: Open-World with V-JEPA~2}
\label{sec:trackb}

Track~B (\emph{Open World}) allows pretrained backbones and external
data. The highest-impact lever in this regime is therefore the choice of
backbone itself: a strong, self-supervised \emph{video} encoder will
dominate any per-frame ResNet trained from scratch. We use
\textbf{V-JEPA~2} \cite{Assran2025VJEPA2}, Meta~FAIR's 2025 video JEPA
that reports state-of-the-art results on motion understanding benchmarks
and, of direct relevance here, the strongest published Something-Something
v2 numbers among open weights: \(77.3\%\) top-1 with the ViT-g/16~384
encoder plus an attentive probe (Tab.~\ref{tab:vjepa2-bench}).

\begin{table}[h]
\centering\small
\begin{tabular}{lcc}
\toprule
\textbf{Benchmark} & \textbf{V-JEPA~2} & \textbf{Previous best} \\
\midrule
SSv2 (probe)     & \textbf{77.3\%} & 69.7\% (InternVideo2-1B) \\
Diving48 (probe) & \textbf{90.2\%} & 86.4\% (InternVideo2-1B) \\
EK100 (probe)    & \textbf{39.7\%} & 27.6\% (PlausiVL) \\
\bottomrule
\end{tabular}
\caption{V-JEPA~2 attentive-probe results reported in
\cite{Assran2025VJEPA2}. SSv2 is the source distribution of the
challenge data set: a frozen V-JEPA~2 encoder almost matches the
previous best \emph{fine-tuned} number.}
\label{tab:vjepa2-bench}
\end{table}

\subsection{Track-B compliance and design choices}

Track~B explicitly allows pretrained weights and external data, so using
the publicly released V-JEPA~2 checkpoints (trained by Meta on
internet-scale video) is permitted. We choose the HuggingFace integration
over the upstream PyTorch~Hub one because the
\code{transformers.AutoModel} path (i) ships a well-tested
\code{VJEPA2Model} class, (ii) caches checkpoints under \code{\$HF\_HOME}
so we can swap variants by name, and (iii) gives access to every
publicly released V-JEPA~2 variant from a single API:
\begin{center}\small
\begin{tabular}{llcc}
\toprule
\textbf{HF repository} & \textbf{Arch.} & \textbf{Params} & \textbf{Res.} \\
\midrule
\code{facebook/vjepa2-vitl-fpc64-256}        & ViT-L & 0.3\,B & 256 \\
\code{facebook/vjepa2-vith-fpc64-256}        & ViT-H & 0.7\,B & 256 \\
\code{facebook/vjepa2-vitg-fpc64-256}        & ViT-g & 1.0\,B & 256 \\
\code{facebook/vjepa2-vitg-fpc64-384}        & ViT-g & 1.0\,B & 384 \\
\code{facebook/vjepa2-vitl-fpc16-256-ssv2}   & ViT-L & 0.4\,B & 256 \\
\code{facebook/vjepa2-vitg-fpc64-384-ssv2}   & ViT-g & 1.0\,B & 384 \\
\bottomrule
\end{tabular}
\end{center}

The two \code{*-ssv2} variants come with an SSv2-finetuned encoder and
are particularly interesting for our 33-class subset: the features have
already been adapted to the Something-Something distribution, so the
probe converges in a handful of epochs.

\paragraph{Why a frozen probe.} Following the standard V-JEPA
evaluation protocol \cite{Assran2025VJEPA2}, the encoder is frozen and
only the classifier probe is trained. This:
\begin{itemize}[leftmargin=2em]
\item Caps trainable parameters at $\sim$\,1\,M (rather than
      $0.3$--$1$\,B), which makes the run fit comfortably on a single
      GPU with batch size~4--8 even at the 384-resolution ViT-g.
\item Avoids destroying the unsupervised representation by overfitting
      on a 33-class slice of 50\,k clips.
\item Reproduces a published number (77.3~\% on the full SSv2) instead
      of inventing a new training recipe.
\end{itemize}

\subsection{The \code{vjepa2} model}
\label{sec:trackb:model}

A new file \file{src/smth2smth/track\_b/vjepa2.py} defines three pieces.

\paragraph{Trunk.} \code{VJEPA2Probe} composes
\code{AutoModel.from\_pretrained(hf\_repo)} (the V-JEPA~2 encoder) with a
trainable head. The encoder is loaded with
\code{attn\_implementation="sdpa"} so PyTorch's scaled-dot-product kernel
is used (FlashAttention-2 when available; eager fallback otherwise).
When \code{freeze\_backbone=true} and \textbf{no} LoRA (Phase~4.4), the
forward path is
\begin{lstlisting}[style=py]
def forward(self, video_batch: torch.Tensor) -> torch.Tensor:
    if self._encoder_needs_grad:
        tokens = self._encode(video_batch)
    else:
        with torch.no_grad():
            tokens = self._encode(video_batch)
    return self.head(tokens)
\end{lstlisting}
With \code{lora\_enabled=true}, \code{\_encoder\_needs\_grad} is true:
the pretrained weights stay frozen (via PEFT) but LoRA adapter paths
receive gradients, so the \code{torch.no\_grad} block is \emph{not} used
on the encoder forward.

\(\code{video\_batch}\) has the shape \((B, T, C, H, W)\) produced by
\code{VideoFrameDataset} with ImageNet normalisation. That layout
matches the HuggingFace \code{VJEPA2Model.forward} contract exactly
(its \cfgkey{pixel\_values\_videos} is documented as
``\code{(batch\_size, num\_frames, num\_channels, frame\_size,
frame\_size)}''), so no reshape is needed.

\paragraph{Frozen-mode invariants.} \code{VJEPA2Probe.train(mode)} keeps
the backbone in \code{eval()} when it is \emph{fully} frozen (no LoRA), so
dropout inside the encoder stays off:

\begin{lstlisting}[style=py]
def train(self, mode: bool = True) -> VJEPA2Probe:
    super().train(mode)
    if self.freeze_backbone and not self.lora_enabled:
        self.backbone.eval()
    return self
\end{lstlisting}

When \code{lora\_enabled=true}, the PEFT-wrapped backbone follows the
parent \code{train()/eval()} so LoRA dropout behaves as intended during
optimisation.

A unit test (\code{test\_train\_keeps\_frozen\_backbone\_in\_eval\_mode})
asserts that, without LoRA, after \code{model.train()} the head is in
train mode but the backbone remains in eval mode. A second test
(\code{test\_frozen\_forward\_does\_not\_accumulate\_backbone\_grads})
asserts that, after backprop with a fully frozen trunk, every backbone
parameter has \code{p.grad is None} while every head parameter has
non-trivial gradient mass (LoRA-enabled runs instead accumulate adapter
gradients only).

\paragraph{Head.} Three head types are offered.
\begin{description}[leftmargin=1.5em, style=nextline]
\item[\code{attentive} (default).]
   The V-JEPA paper's attentive probe, \emph{generalised} to \(Q\ge 1\)
   learnable query tokens \(q \in \mathbb{R}^{Q\times D}\). Each query
   cross-attends over the \((B, N, D)\) token sequence; the \(Q\) pooled
   vectors are \textbf{mean-pooled} across queries, then
   \(\mathrm{LayerNorm}\) and a linear classifier are applied:
   \[
   Z \;=\; \mathrm{MHA}(q,\, K{=}V{=}\Phi)\in\mathbb{R}^{B\times Q\times D},
   \qquad
   z \;=\; \tfrac{1}{Q}\sum_{i=1}^{Q} Z_{:,i,:},
   \qquad
   \widehat{y} = Wz + b.
   \]
   The config key \cfgkey{head\_num\_queries} controls \(Q\); the default
   \(Q{=}1\) recovers the original single-query recipe. Larger \(Q\) adds
   probe capacity for compositional verbs without widening the backbone.
\item[\code{linear} / \code{mean\_linear}.]
   Mean-pool the token sequence and project with a single Linear.
   Fewer parameters, marginally faster; useful as an ablation baseline.
\end{description}

\paragraph{Builder.} A \code{@register\_model("vjepa2")} decorator
exposes the model under the same registry the rest of the project uses:

\begin{lstlisting}[style=py]
@register_model("vjepa2")
def build_vjepa2(cfg: DictConfig) -> nn.Module:
    return VJEPA2Probe(
        num_classes=int(cfg.model.num_classes),
        hf_repo=str(cfg.model.get("hf_repo", DEFAULT_HF_REPO)),
        head_type=str(cfg.model.get("head_type", "attentive")),
        head_num_heads=int(cfg.model.get("head_num_heads", 8)),
        head_num_queries=int(cfg.model.get("head_num_queries", 1)),
        head_dropout=float(cfg.model.get("head_dropout", 0.0)),
        freeze_backbone=bool(cfg.model.get("freeze_backbone", True)),
        lora_enabled=bool(cfg.model.get("lora_enabled", False)),
        lora_r=int(cfg.model.get("lora_r", 8)),
        lora_alpha=int(cfg.model.get("lora_alpha", 16)),
        lora_dropout=float(cfg.model.get("lora_dropout", 0.05)),
        lora_target_modules=...,
        attn_implementation=str(cfg.model.get("attn_implementation", "sdpa")),
    )
\end{lstlisting}
The ellipsis stands for optional \code{lora\_target\_modules} parsed from
Hydra (default \code{null} \(\Rightarrow\) HF V-JEPA~2
\code{query}, \code{key}, \code{value}, \code{proj}).

\subsection{Registration hook}

\code{shared/models/\_\_init\_\_.py} imports \code{smth2smth.track\_b} at
module load time so that the \code{@register\_model("vjepa2")} decorator
runs before \code{build\_model(cfg)} is ever called. This is the only
cross-package import between \code{shared/} and \code{track\_b/}; the
opposite direction is empty (\code{track\_b/vjepa2.py} only imports from
\code{shared.models.registry}), so there is no circularity:
\begin{lstlisting}[style=py]
% src/smth2smth/shared/models/__init__.py
from smth2smth import track_b as _track_b  % noqa: F401
\end{lstlisting}

A unit test (\code{test\_vjepa2\_is\_registered}) asserts that
\code{"vjepa2"} is in \code{list\_registered\_models()} after a clean
import. A second test (\code{test\_build\_model\_dispatches\_via\_registry})
goes one step further and runs \code{build\_model(cfg)} with a mocked
HuggingFace backbone (\code{sys.modules["transformers"]} is replaced
with a stub) to check that the dispatch works end-to-end through the
public API.

\subsection{The bundled experiment configuration}

\file{configs/experiment/track\_b\_vjepa2.yaml} composes the V-JEPA~2
model with sensible Track-B training defaults:

\begin{lstlisting}[style=yaml]
defaults:
  - override /model: vjepa2

model:
  pretrained: true                     # ImageNet normalisation
  hf_repo: facebook/vjepa2-vitl-fpc64-256
  head_type: attentive
  head_num_heads: 8
  freeze_backbone: true

dataset:
  num_frames: 16                       # must be even (tubelet_size=2)
  image_size: 256                      # match the backbone's crop_size
  use_official_val: true

training:
  batch_size: 4
  epochs: 30
  lr: 0.001
  optimizer: adamw
  weight_decay: 0.05
  scheduler_cosine: true
  warmup_epochs: 2
  label_smoothing: 0.1
  amp: true
  ema_enabled: true
  ema_decay: 0.999
  tta: true
  tta_flip: true
\end{lstlisting}

Three settings are non-negotiable:
\begin{enumerate}[leftmargin=2em]
\item \cfgkey{model.pretrained: true}.\\
      \code{build\_transforms} reads this and applies the ImageNet mean
      and std, which is what V-JEPA~2 was trained with. Setting it to
      \code{false} would normalise with \((0.5, 0.5, 0.5)\) and silently
      give wrong inputs to the encoder.
\item \cfgkey{dataset.image\_size: 256} (or 384 for the
      \code{vitg-384} variants).\\
      The V-JEPA~2 encoders are tied to a fixed crop size by their
      learned positional embeddings; resizing inputs to anything else is
      not supported by the public weights.
\item \cfgkey{dataset.num\_frames} must be even.\\
      The V-JEPA~2 tubelet size is~2; an odd \(T\) is not handled by the
      encoder's temporal patchifier and raises a runtime error.
\end{enumerate}

The remaining settings (AMP, EMA, TTA, label smoothing, cosine schedule
with 2-epoch warmup) carry over from the Phase-2 recipe with no
changes, because the supervised \code{train.py} is generic over
\code{cfg.model.name}.

\subsection{Hydra symbol glossary (training runs)}
\label{sec:trackb:glossary}

Every long training job we care about is summarised in
Tab.~\ref{tab:trackb-runlog}.  Hyperparameters are written in
\textbf{compact form} using the symbols below; the full English meaning
lives here once, not repeated in each run row.

\begin{table}[ht]
\centering\footnotesize
\begin{tabular}{@{}r p{0.32\linewidth} p{0.52\linewidth}@{}}
\toprule
\textbf{Sym.} & \textbf{Hydra key(s)} & \textbf{Meaning} \\
\midrule
\texttt{expt} & \texttt{experiment=\dots} & Preset YAML under
\file{configs/experiment/} (composed defaults). \\
\texttt{repo} & \texttt{model.hf\_repo} & HuggingFace V-JEPA~2 checkpoint id. \\
\texttt{T} & \texttt{dataset.num\_frames} & Frames per clip (must be even). \\
\texttt{Sz} & \texttt{dataset.image\_size} & Square crop side (256 or 384). \\
\texttt{bs} & \texttt{training.batch\_size} & Clips per optimisation step. \\
\texttt{E} & \texttt{training.epochs} & Total epoch budget for the run. \\
\texttt{lr} & \texttt{training.lr} & Base learning rate (AdamW). \\
\texttt{wd} & \texttt{training.weight\_decay} & AdamW weight decay. \\
\texttt{wu} & \texttt{training.warmup\_epochs} & Linear LR warm-up length. \\
\texttt{cos} & \texttt{training.scheduler\_cosine} & Cosine anneal after warm-up. \\
\texttt{ls} & \texttt{training.label\_smoothing} & Soft-label smoothing for CE. \\
\texttt{amp} & \texttt{training.amp} & Automatic mixed precision. \\
\texttt{ema} & \texttt{training.ema\_enabled} & Exponential moving average
of weights (0/1); decay in \texttt{ema\_decay}. \\
\texttt{ev} & \texttt{training.eval\_every\_n\_epochs} & Validate every
\(n\) epochs. \\
\texttt{mz} & \texttt{dataset.max\_samples\_per\_class} & Cap clips per class
(\texttt{null} = use all). \\
\texttt{aug} & \texttt{defaults} / augment group & Spatial/temporal augment
preset (e.g. \texttt{vjepa2\_heavy}). \\
\texttt{LoRA} & \texttt{model.lora\_enabled}, \texttt{lora\_r},
\texttt{lora\_alpha} & PEFT adapters on attention projections; pretrained
weights stay frozen. \\
\texttt{ckpt} & \texttt{training.checkpoint\_path} & Primary weights file. \\
\texttt{rs} & \texttt{training.resume\_from} & Optional optimiser/epoch
resume path. \\
\bottomrule
\end{tabular}
\caption{Symbols used in Tab.~\ref{tab:trackb-runlog}.  Extra CLI overrides
are appended verbatim after a semicolon.}
\label{tab:trackb-glossary}
\end{table}

\subsection{Track B training run logbook}
\label{sec:trackb:runlog}

Statuses: \textbf{RUNNING} (PID alive or log still progressing),
\textbf{DONE} (trainer printed \texttt{Done.}), \textbf{FAILED} (traceback
or non-zero exit).  After any merge or local pull, refresh this table
from \texttt{logs/} and PID files.

\begin{table}[ht]
\caption{Chronological log of substantial Track-B training jobs.  Update
status and metrics when processes finish; add new rows for every fresh
\texttt{nohup} launch.}
\label{tab:trackb-runlog}
\begin{minipage}{\linewidth}\small
\begin{description}[leftmargin=1.5em, style=nextline]
\item[\textbf{Run 1} (DONE)]
  Log: \texttt{logs/track\_b\_vitg384\_heavy\_aug\_20260512\_0159.log}.
  Preset: \texttt{track\_b\_vjepa2\_vitg384\_heavy\_aug}.
  Params (Tab.~\ref{tab:trackb-glossary}): \texttt{repo}=vitg-fpc64-384-ssv2;
  \texttt{T}=16; \texttt{Sz}=384; \texttt{bs}=4; \texttt{E}=10; \texttt{lr}=3e-4;
  \texttt{wd}=0.05; \texttt{wu}=1; \texttt{mz}=100; \texttt{ev}=5; \texttt{ema}=0.
  Best val top-1: \textbf{53.68\,\%}.
\item[\textbf{Run 2} (DONE)]
  Log: \texttt{logs/track\_b\_vitg384\_continue\_20260512\_1322.log};
  PID: \texttt{logs/track\_b\_vitg384\_continue.pid}.
  Preset: \texttt{track\_b\_vjepa2\_vitg384\_continue}.
  Params: inherits Run~1 + \texttt{rs}=vitg384\_heavy\_aug.last.pt;
  \texttt{E}=20 (epochs 11--20); \texttt{wu}=0; \texttt{ev}=1; \texttt{cos}=0.
  Best val top-1: \textbf{53.94\,\%}.
\item[\textbf{Run 3} (RUNNING)]
  Log: \texttt{logs/track\_b\_vitl\_30e\_warmup\_lora\_20260514\_1536.log};
  PID: \texttt{logs/track\_b\_vitl\_30e\_warmup.pid}.
  Preset: \texttt{track\_b\_vjepa2\_vitl\_30e\_warmup}.
  Params: \texttt{repo}=vitl-fpc64-256; \texttt{T}=16; \texttt{Sz}=256; \texttt{bs}=4;
  \texttt{E}=30; \texttt{wu}=3; \texttt{ev}=1; \texttt{ema}=0; LoRA \texttt{r}=8, $\alpha$=16.
  Interim (epoch~1) val top-1: \textbf{60.21\,\%}.
\end{description}
\end{minipage}
\end{table}

\subsection{Running V-JEPA 2 on Track B}

The sole additional dependency is \code{transformers}~$\geq 4.45$, the
release that ships \code{VJEPA2Model}; training, evaluation and
submission all go through the existing \code{smth2smth.pipelines}
entrypoints with \code{track=b} and the appropriate experiment config.
The encoder is downloaded from HuggingFace Hub on first use
($\approx$\,1.2\,GiB for \code{vitl-fpc64-256}) and cached under
\code{\$HF\_HOME}; subsequent runs reuse the cache.

% =====================================================================
% Phase 4.1: zero-shot evaluation and submission with the SSv2-finetuned
% V-JEPA 2 checkpoint -- run \emph{without any training of our own}.
% =====================================================================

\subsection{Zero-shot Evaluation with the SSv2 Head}
\label{sec:trackb:zero_shot}

The SSv2-finetuned checkpoint
\code{facebook/vjepa2-vitl-fpc16-256-ssv2} ships a full 174-class
classifier head trained by Meta on the full SSv2 distribution. Our
33-class subset is a strict subset of those 174 classes, so we can
evaluate -- and submit -- with no training of our own. Two small
artifacts implement this in the codebase:

\begin{itemize}[leftmargin=2em]
\item \file{src/smth2smth/track\_b/zero\_shot.py} -- reusable helpers:
   \code{normalize\_class\_name}, \code{build\_label\_mapping} (with a
   token-aligned unique-prefix fallback that recovers the truncated
   folder \code{015\_Pretending\_to\_pour\_something\_out\_of\_something\_but\_something\_})
   and \code{VideoFramesDataset} that reads 16 evenly-spaced frames per
   clip, resizes to 256, ImageNet-normalises, and stacks to \((T,3,H,W)\).
\item \file{scripts/eval\_vjepa2\_ssv2\_zero\_shot.py} -- CLI that loads
   \code{VJEPA2ForVideoClassification}, slices the 174-d logits to the
   matched 32 classes, and reports top-1, top-5, and a per-class
   accuracy table.
\item \file{scripts/submit\_vjepa2\_ssv2\_zero\_shot.py} -- same
   forward, but consumes \code{data/test/} and writes
   \code{submissions/track\_b\_vjepa2\_vitl\_fpc16\_256\_ssv2\_zero\_shot.csv}
   in the official two-column format; the file is re-validated with
   \code{validate\_submission\_csv} before exit. An optional
   \code{--tta-flip} flag mirrors the horizontal-flip TTA + auto
   left/right class-pair remap that \code{pipelines/submit.py} performs
   for trained checkpoints.
\end{itemize}

\paragraph{Result.} On the official \code{data/val/} split (6\,745
clips, 32 of 33 classes -- class 27 has no validation data) the
zero-shot top-1 is \(\mathbf{44.0\%}\) with top-5 at
\(\mathbf{74.5\%}\), in 373\,s on a single NVIDIA RTX~A4000
(18 vid/s). The public Kaggle leaderboard score on the
\(6\,913\)-video test set is \(\mathbf{45\%}\) top-1 -- already
\(\sim 8\)\,pp ahead of the Track-A from-scratch baseline (37.2\%).
The result is documented in detail in \code{notebooks/06\_vjepa2\_feature\_geometry.ipynb}
together with PCA and t-SNE visualisations of the encoder's
pre-classifier features, which confirm that many of the 32 classes form
visually well-separated clusters in 2D even before any training of our
own --- predicting that a thin probe should comfortably beat the
zero-shot baseline.

% =====================================================================
% Phase 4.2: heavy data augmentation for the V-JEPA 2 probe.
% Mirrors the Track-A RandAugment-T + FrameMixUp recipe (Kim et al.
% 2020) and reuses the existing augmentation infrastructure that was
% built in Phase 2.
% =====================================================================

\subsection{Data Augmentation for the V-JEPA~2 Probe}
\label{sec:trackb:heavy_aug}

The default \file{track\_b\_vjepa2.yaml} above only composes the model
group; the augmentation group falls back to \code{augment=none}
(resize + horizontal flip). That is a deliberate choice for the
zero-shot evaluation, but for any training run we want to inject as
much randomness as possible: the V-JEPA~2 encoder is frozen (only
\(\sim\,4\,\)M probe parameters receive gradients) so the probe will
overfit our 33-class slice in a handful of epochs unless we vary the
features it sees per epoch. Heavy augmentation does exactly that.

\paragraph{Mirroring Track A.}
The Track-A side of the project (\code{track\_a\_randaugment\_t.yaml})
already implements Kim~et~al.~2020's best recipe (Tab.~7): RandAugment-T
(\code{mode: temporal\_plus}) layered with FrameMixUp. The same
infrastructure transfers to Track~B without code changes -- both
\code{shared/data/transforms.py} (per-frame augmentations) and
\code{shared/engine/trainer.py} (clip-level VideoMix) are agnostic to
the model name and the track. The only thing missing is the Track-B
experiment that composes these knobs in the right place.

\paragraph{What \code{augment=vjepa2\_heavy} stacks.}
\file{configs/augment/vjepa2\_heavy.yaml} (kept from earlier
exploration) chains, in order, with all randomness \emph{synchronised
across the frames of a clip} so that motion direction is preserved:
\begin{enumerate}[leftmargin=2em]
\item \textbf{RandAugment-T} (Kim et~al. 2020, Sec.~3.1+4.2) -- 14-op
   palette, \(N{=}2\), \(M{=}9\), \code{magnitude\_max=30},
   \code{mode: temporal\_plus} which is the strongest setting in
   Tab.~2: a magnitude \(M\) is drawn per clip and the per-frame
   magnitude varies linearly from \(M-\delta\) to \(M+\delta\) with
   \(\delta\sim U(0, 0.5M)\). This injects temporally-coherent
   geometric/photometric noise that mimics camera motion and lighting
   drift.
\item \textbf{Random crop} 256\,+\,32~\(\to\)~256, mild zoom jitter
   that prevents the SSv2 prior from latching onto absolute pixel
   positions.
\item \textbf{Random horizontal flip}. The class-pair remap for
   ``pulling left~\(\to\)~right'' vs ``right~\(\to\)~left'' lives in
   the submission TTA path, not at train-time -- the flip here just
   adds variety.
\item \textbf{Color jitter} (\code{brightness=0.3}, \code{contrast=0.3},
   \code{saturation=0.3}, \code{hue=0.05}) on top of RandAugment.
\end{enumerate}

\paragraph{Clip-level VideoMix.}
On top of the spatial augmentation, the Track-B heavy-aug experiment
also enables one of the
spatio-temporal mix strategies plumbed in \code{trainer.py}: we
default to \code{videomix\_mode: cube\_cutmix} with
\code{videomix\_alpha: 1.0} and \code{videomix\_prob: 0.5}. CubeCutMix
\cite{Yun2019CutMix} samples a 3-D spatio-temporal sub-volume from a
second clip in the same batch and pastes it onto the current clip;
labels are mixed in proportion to the sub-volume's relative volume.
This is the regularizer of choice for motion-sensitive datasets
(Kim~et~al.~2020 Tab.~7), where naive MixUp blurs motion cues across
the entire clip.

\paragraph{Slim validation loop.}
\code{pipelines/train.py} historically runs two full validation passes
per epoch (live~+~EMA), each at 18\,vid/s on the A4000 \(\to\) 12
minutes for the 6\,745-clip val split. With heavy aug + 15 epochs that
adds up to \(\sim\)\,3\,h of pure-eval wall time, dwarfing the actual
training. Two new Hydra knobs in \code{configs/train/default.yaml}
gate this without breaking any earlier experiment:
\begin{lstlisting}[style=yaml]
training:
  eval_every_n_epochs: 1        # default: eval every epoch (legacy)
  eval_ema: true                # default: also evaluate the EMA snapshot
\end{lstlisting}
With \code{eval\_every\_n\_epochs: 3} and \code{eval\_ema: false} we
shrink the per-epoch overhead from 12\,min to 4\,min averaged over the
schedule, which is exactly the optimisation the per-epoch timing run
(Sec.~\ref{sec:trackb:heavy_aug}) flagged as the dominant cost.

\paragraph{The bundled \code{track\_b\_vjepa2\_heavy\_aug} experiment.}
\texttt{configs/\allowbreak experiment/\allowbreak track\_b\_vjepa2\_heavy\_aug.yaml}
composes the V-JEPA~2 SSv2 backbone with the heavy aug, CubeCutMix, the slim
validation loop, and a small-subset training schedule that fits a
single run in about two hours on the A4000:
\begin{lstlisting}[style=yaml]
defaults:
  - override /model: vjepa2
  - override /augment: vjepa2_heavy

model:
  pretrained: true
  hf_repo: facebook/vjepa2-vitl-fpc16-256-ssv2
  head_type: attentive
  head_num_heads: 8
  head_dropout: 0.1
  freeze_backbone: true

dataset:
  num_frames: 16
  image_size: 256
  use_official_val: true
  max_samples_per_class: 50

training:
  batch_size: 12
  epochs: 15
  lr: 0.0005
  optimizer: adamw
  weight_decay: 0.05
  scheduler_cosine: true
  warmup_epochs: 2
  min_lr: 0.00001
  label_smoothing: 0.1
  amp: true
  ema_enabled: true
  ema_decay: 0.999
  eval_every_n_epochs: 3
  eval_ema: false
  videomix_mode: cube_cutmix
  videomix_alpha: 1.0
  videomix_prob: 0.5
  tta: true
  tta_flip: true
\end{lstlisting}

\paragraph{Why this is the "smallest" V-JEPA 2.}
The smallest V-JEPA~2 SSv2-finetuned checkpoint available on
HuggingFace Hub is \code{vitl-fpc16-256-ssv2} (0.4\,B parameters).
V-JEPA~2.1 ViT-B (80\,M, distilled from ViT-G) was announced by Meta
in March~2026 \cite{MurLabadia2026VJEPA21} but is not yet uploaded
to the Hub at the time of writing; a tracking PR on
\code{transformers} adds inference support, and once weights are
public swapping \code{model.hf\_repo} is the only change required.
For now ViT-L 256 is both the smallest SSv2-finetuned encoder and
also the encoder behind the 44\,\% zero-shot floor, so we keep it as
the iteration baseline.

\paragraph{Why we did not add extra positional encoding.}
V-JEPA~2 already includes a learned 3-D positional embedding per token
(\(T_{\text{tubelet}}{=}2\) along time \(\times\) \(H/16\,\times\,W/16\)
in space, 2\,048 tokens for 16 frames at 256\,px). The attentive
probe's cross-attention reads from this positionally-aware token
sequence with a learnable query, so the global pooling already has
full positional context. Stacking another positional embedding on top
would be redundant. The ceiling we hit at the zero-shot baseline is
\emph{informational}, not positional: \code{data/train/} only contains
four unique frames per video, repeated 4\(\times\) to fill 16
tubelet-aligned positions, so motion-direction-only classes (closing
vs opening, moving-up vs moving-down) lose information at the source.
Mitigations on this front (frame interpolation, optical flow channel)
are deferred to a later phase.

\subsection{Overnight ViT-g/14 384 Run}
\label{sec:trackb-vitg384}

\paragraph{Motivation.}
The Phase~4.2 ViT-L 256 run topped out at \(53.09\,\%\) val top-1
(checkpoint \code{checkpoints/track\_b/vjepa2\_heavy\_aug.pt}, 15
epochs, 200 clips/class, EMA off, no augmentation). That is only
\(+8.7\) percentage points over the \(44\,\%\) zero-shot floor and
\emph{below} the \(>70\,\%\) anchor reported by peer submissions. The
bottleneck is not augmentation strength but \emph{backbone capacity}:
a 0.4\,B-parameter ViT-L probed on 6\,400 clips is now competing with
1\,B-parameter encoders trained on the same task. The Phase~4.3 run
swaps to the largest SSv2-finetuned V-JEPA~2 available on the Hub --
\code{facebook/vjepa2-vitg-fpc64-384-ssv2} (1.0\,B parameters, 384\,px
input, \(77.3\,\%\) SSv2 in the paper) -- and combines it with the
Phase~4.2 augmentation recipe. To keep wall time inside an overnight
slot we shrink the training subset accordingly.

\paragraph{Memory and throughput smoke check.}
On the lab A4000 (16\,GiB total, 12.2\,GiB free at launch) a forward
pass with \code{torch.float16} and \code{skip\_predictor=True} measures
2.27\,GiB resident after weight load and the following inference-only
costs:
\begin{itemize}
  \item batch 1, 16 frames @ 384\,px: 882\,ms/clip (\(\approx 1.13\)
        vid/s), 2.49\,GiB peak.
  \item batch 2: 1\,548\,ms/step (\(1.29\) vid/s), 2.70\,GiB peak.
  \item batch 4: 3\,148\,ms/step (\(1.27\) vid/s), 3.10\,GiB peak.
  \item batch 6: 4\,553\,ms/step (\(1.32\) vid/s), 3.52\,GiB peak.
\end{itemize}
Throughput plateaus at \(\sim\!1.3\) vid/s -- the encoder is
memory-bandwidth-bound, so larger batches mostly reduce host-side
overhead. We pick batch 4 to keep the optimiser stable while staying
well below the 12\,GiB free budget.

\paragraph{Recipe.}
The experiment lives at
\code{configs/experiment/track\_b\_vjepa2\_vitg384\_heavy\_aug.yaml}
and overrides Phase~4.2 with the following changes:
\begin{itemize}
  \item \code{model.hf\_repo: facebook/vjepa2-vitg-fpc64-384-ssv2}
        (frozen, \(\sim\!1.0\,\)B params, \(\sim\!8.0\,\)M trainable
        in the attentive probe).
  \item \code{dataset.num\_frames: 16}, \code{image\_size: 384}.
  \item \code{dataset.max\_samples\_per\_class: 100} (\(100\!\times\!32
        = 3\,200\) train clips, balanced).
  \item \code{training.batch\_size: 4}, \code{epochs: 10}, cosine LR
        with 1 warm-up epoch, \(\text{lr}{=}3\!\times\!10^{-4}\),
        \(\text{wd}{=}0.05\), label smoothing 0.1.
  \item \code{augment=vjepa2\_heavy} (RandAugment-T temporal\_plus +
        random crop + flip + colour jitter) and clip-level CubeCutMix
        (\(\alpha{=}1.0\), \(p{=}0.5\)).
  \item Slim validation: \code{eval\_every\_n\_epochs: 5} (only
        2 val passes total, epochs 5 and 10), \code{ema\_enabled:
        false} so no \(\sim\!4\,\)GiB deepcopy.
  \item TTA flip + class-pair remap at submission time.
\end{itemize}

\paragraph{Resume-safe checkpointing.}
A long run that crashes between val checkpoints can lose hours of
work: the previous trainer only wrote \code{best\_model.pt} on
validation improvements. We added two new \code{training.\,*} knobs
(see \code{configs/train/default.yaml}):
\begin{itemize}
  \item \code{save\_last\_checkpoint: true} -- end of every epoch the
        live model state, optimiser state, cosine scheduler state and
        AMP \code{GradScaler} state are written to a companion file.
  \item \code{last\_checkpoint\_path: null} -- by default derived from
        \code{checkpoint\_path} by inserting \code{.last} before the
        extension (\code{vitg384\_heavy\_aug.pt} \(\to\)
        \code{vitg384\_heavy\_aug.last.pt}).
\end{itemize}
On a crash, restart with
\begin{verbatim}
PYTHONPATH=src python -m smth2smth.pipelines.train \
  track=b experiment=track_b_vjepa2_vitg384_heavy_aug \
  training.checkpoint_path=checkpoints/track_b/vitg384_heavy_aug.pt \
  training.resume_from=checkpoints/track_b/vitg384_heavy_aug.last.pt
\end{verbatim}
The trainer reads back the saved epoch counter, restores
optimiser/scheduler/scaler state, and resumes at the next epoch
boundary. The existing
\code{tests/pipelines/test\_last\_checkpoint.py} unit-tests pin all
three contract points (write-every-epoch, opt-out, override path).

\paragraph{Launch command.}
Detached, with a timestamped log file and a PID file for monitoring:
\begin{verbatim}
PYTHONPATH=src nohup .venv/bin/python -u \
  -m smth2smth.pipelines.train \
  track=b experiment=track_b_vjepa2_vitg384_heavy_aug \
  training.checkpoint_path=checkpoints/track_b/vitg384_heavy_aug.pt \
  hydra.run.dir=outputs/track_b_vitg384_heavy_aug_<stamp> \
  > logs/track_b_vitg384_heavy_aug_<stamp>.log 2>&1 < /dev/null &
echo $! > logs/track_b_vitg384_heavy_aug.pid
\end{verbatim}
\code{nohup} keeps the run alive after the shell disconnects;
\code{-u} forces unbuffered stdout so the log file reflects progress
in real time; \code{< /dev/null} stops the worker from blocking on
\code{stdin}.

\paragraph{Wall-time projection (overnight slot).}
With \(800\) training steps per epoch (\(3\,200\) clips / batch 4)
and the smoke-measured \(\sim\!1.3\) vid/s effective training
throughput (encoder forward dominates; backward only flows through
the 8\,M-param probe), one epoch costs \(\approx 41\,\)min, observed
to be \(\approx 45\,\)min in the actual log including data-loader
overhead. The full \code{data/val/} pass costs \(\sim\!78\,\)min at
the same throughput. Total budget: \(10\times 45 + 2 \times 78
\approx 606\,\)min, i.e. \(\sim\!10\,\)h -- fits a
01:59\,\(\to\)\,12:00 window.

\paragraph{Outcome.}
Finished ViT-g~384 heavy-aug jobs are summarised in
Tab.~\ref{tab:trackb-runlog} (rows~1--2) with compact notation from
Tab.~\ref{tab:trackb-glossary}.  Phase~4.3 supersedes Phase~4.2 as the
flagship Track-B recipe for large backbones; the submission CSV is
produced via \code{python -m smth2smth.pipelines.submit} pointing at
\code{checkpoints/track\_b/vitg384\_heavy\_aug.pt}.

\subsection{Multi-query Probe and PEFT LoRA}
\label{sec:trackb-peft}

\paragraph{Motivation.}
The gap between course val (\(\sim\!54\%\) with ViT-g~384) and the
paper's \(77.3\,\%\) SSv2 benchmark is only partly model-size: the
frozen encoder cannot adapt its internal features to the 33-class
slice or to domain shift, while a single-query attentive probe may
under-represent compositional verbs. Phase~4.4 adds two orthogonal
extensions implemented in \file{src/smth2smth/track\_b/vjepa2.py} and
exposed through \file{configs/model/vjepa2.yaml}:

\begin{enumerate}[leftmargin=2em]
\item \textbf{Multi-query attentive head} ---
      config key \cfgkey{head\_num\_queries} (default \code{1}).
      \(Q\) learnable queries each run cross-attention over the full
      token grid; outputs are mean-pooled across \(Q\) before
      \code{LayerNorm} and the classifier. This increases probe capacity
      without touching the ViT trunk.
\item \textbf{PEFT LoRA} --- optional low-rank adapters (LoRA, Hu et al.,
      ICLR~2022) on attention linear projections
      (\code{query}, \code{key}, \code{value}, \code{proj} by default for
      the HuggingFace \code{VJEPA2Model} encoder; configurable via
      \cfgkey{lora\_target\_modules}).
      The \texttt{peft} package wraps the HuggingFace encoder with
      \code{get\_peft\_model}; pretrained weights stay frozen while
      adapter matrices train. Forward passes run \emph{with} autograd
      through the adapters (\code{\_encoder\_needs\_grad=true}), unlike
      the pure frozen trunk. Full encoder fine-tuning remains available
      via \code{freeze\_backbone=false} with \code{lora\_enabled=false}
      (not combinable with LoRA in the current implementation).
\end{enumerate}

\paragraph{Dependencies.}
The project declares \texttt{peft>=0.13} in \file{pyproject.toml}.
LoRA is inactive unless \code{model.lora\_enabled=true}.

\paragraph{Experiment preset.}
\file{configs/experiment/track\_b\_vjepa2\_peft.yaml} composes the Phase~4.3
ViT-g~384 heavy-aug recipe with \code{head\_num\_queries: 4},
\code{lora\_enabled: true}, \code{lora\_r: 8}, \code{lora\_alpha: 16},
and a slightly reduced learning rate (\code{training.lr: 2e-4}) so
adapter updates remain stable. Checkpoints are written to
\code{checkpoints/track\_b/vitg384\_peft\_mq.pt} by default.

\paragraph{Caveats.}
If PEFT cannot match \code{lora\_target\_modules} to the encoder's
module names, construction raises a clear \code{RuntimeError}; adjust
the list for the specific \code{hf\_repo}. QLoRA (4-bit base weights) is
\emph{not} implemented here---only BF16/FP16 training as before. Dense
multi-crop / multi-clip test-time ensembling remains future work.
